%------------------------------------------------------------------------------%
\begin{question}
  Dada a equação da circunferência
  \[
    (x-2)^2 + (y+3)^2 = 25
  \]
  \vspace{-1.5\baselineskip}
  \begin{enumerate}
    \item encontre seu raio e o centro
    \item escreva a equação matricial transladada
    \item esboce a curva
  \end{enumerate}
\end{question}

%------------------------------------------------------------------------------%
\begin{resolution}{Raio de Centro}
  Comparando
  \[
    (x-2)^2 + (y+3)^2 = 25
  \]
  com a equação canônica transladada
  \[
    (x-h)^2 + (y-k)^2 = 25
  \]

  Identificamos

  \quad o centro:
  \(
    C
    = (h, k)
    = (2, -3)
  \)

  \quad o raio:
  \(
    r
    = \sqrt{25}
    = 5
  \)
\end{resolution}

%------------------------------------------------------------------------------%
\begin{resolution}{Equação Matricial}
  Equação Matricial
  \[
    V^t I V = r^2
  \]
  com
  \[
    I = \begin{bmatrix}
      1 & 0 \\
      0 & 1
    \end{bmatrix}
  \]
  \[
    V
    = \begin{bmatrix}
      x - h  \\
      y - k
    \end{bmatrix}
    = \begin{bmatrix}[c]
      x - 2  \\
      y - (-3)
    \end{bmatrix}
    = \begin{bmatrix}
      x - 2  \\
      y + 3
    \end{bmatrix}
  \]

  \[
    r = 5
  \]
\end{resolution}

%------------------------------------------------------------------------------%
\begin{resolution}{Esboço}
\end{resolution}

%------------------------------------------------------------------------------%
\endinput
